%% ===========================================================================
\FloatBarrier
\section{Parameter identification and limitations}
\label{sec:calibration}
%% ===========================================================================

\subsection{Measurement routes}
\label{sec:calib_a_mu}

The window size $a$, fractional order $\alpha$, relaxation time $\tau$ and
flexoelectric coefficient $\mu$ require identification rather than a
constitutive-table lookup. Size-dependent compliance or frequency identifies
$M_{22}$; directional measurements identify $M_{11}/M_{22}$ and the window
orientation, as in characteristic-length identification~\cite{fem1}.
Ultrasonic dispersion or a clamped-edge strain profile provides an independent
length estimate. Microstructure images and two-point correlations can determine
the window anisotropy and test whether its second moment is an adequate
description. These are possible calibration routes, not measurements performed
in this study.

The flexoelectric coefficient follows from voltage or charge at a measured
curvature $\kappa_c$; for the bending configuration of
Fig.~\ref{fig:geometry}(a),
\begin{equation}
  V = \frac{\mu\,\kappa_c\,h}{\kappa_{33}},
  \label{eq:muV}
\end{equation}
where $\kappa_{33}=4.9583\times10^{-10}$\,\si{\farad\per\meter}.
Equation~\eqref{eq:muV} identifies the magnitude; the sign depends on crystal
symmetry. The assumed $\mu=10^{-6}$\,\si{\coulomb\per\meter} sets absolute
bending couplings, which scale as $\mu^2$.

\subsection{Fractional order and relaxation time}
\label{sec:calib_alpha_tau}

Dielectric spectroscopy identifies $(\alpha,\tau)$ by jointly fitting the
modulus and loss tangent of \eqref{eq:kvfreq},
\begin{equation}
  \tan\delta=\frac{(\omega\tau)^{\alpha}\sin(\pi\alpha/2)}
     {1+(\omega\tau)^{\alpha}\cos(\pi\alpha/2)}.
  \label{eq:tandelta}
\end{equation}
Both are monotone in $\omega\tau$, so a single loss peak does not identify
$\tau$. Ring-down measurements provide a complementary route; for small loss,
\begin{equation}
  \zeta=\frac{(\omega\tau)^{\alpha}\sin(\pi\alpha/2)}{2}
  \quad\Longrightarrow\quad
  \delta\approx\pi(\omega\tau)^{\alpha}\sin(\pi\alpha/2).
  \label{eq:zetasmall}
\end{equation}
At $\alpha=1$, \eqref{eq:zetasmall} gives $\delta=\pi\omega\tau$, agreeing
with \eqref{eq:zeta} to the $1.2\,\%$ small-loss approximation. One decrement
constrains only a product; identifying both parameters requires multiple
frequencies or an independently known order. The fitted law represents a
broadened single-power-law response, not two separated loss mechanisms.

\subsection{Validity window}
\label{sec:limit_window}

The geometry and parameter ranges are those of Table~\ref{tab:params} and the
individual sweeps. The damping interpretation is restricted to
$\omega\tau\le1$; outside that band the map extrapolates the fitted law.
The second-moment truncation uses $a/L\le0.2$ as an assumed operating range,
not a validated truncation threshold: no experiment or exact averaging
comparison here establishes where the approximation fails. The coupled
denominator of \eqref{eq:kstatic} is retained because $k^2=0.5581$ is not
small; the uncoupled measure is used only for the constitutive comparison.

\subsection{Modeling limits}
\label{sec:limitations}

The results use two-dimensional plane stress; plane strain gives $0.5132$
for the constitutive coupling, and neither reduction is a three-dimensional
validation. Density rescales absolute frequencies, while $\mu^2$ rescales
absolute flexoelectric couplings. Memory placement must accompany calibration:
$K_{\rm eff}$ and $K_{uu}$ give a damping-ratio ratio of $0.450$ at $\alpha=1$
for the device of Sec.~\ref{sec:memorystudy}. They are distinct condensed
models, not interchangeable numerical implementations. Thin-strip couplings
are resolved to two significant digits because they subtract nearly equal
eigenvalues. Verification establishes that the implementation solves the
stated model; material-specific prediction requires measured inputs.

\FloatBarrier
